\begin{table}[H]
\centering
\begin{threeparttable}
\caption{Follow-up Survey: Retention, Media Evaluation, and Later-Demand Modules}
\label{tab:followup_survey_selection}
\scriptsize
\setlength{\tabcolsep}{4pt}
\begin{tabularx}{\textwidth}{p{7.1cm} p{3.4cm} X}
\toprule
\textbf{Question or module} & \textbf{Response format} & \textbf{Purpose in analysis} \\
\midrule
Please rate your feelings toward China and toward Western democracies. & Two 0--100 sliders & Persistence in the China minus Western-democracy thermometer gap \\
\addlinespace
Factual questions about details from the earlier reading set. & Twelve multiple-choice items scored 0/1 & Knowledge retention after treatment \\
\addlinespace
How biased are Chinese media when reporting on China? & 0--10 slider & Media evaluation \\
\addlinespace
How biased are Western media when reporting on China? & 0--10 slider & Media evaluation \\
\addlinespace
Would you like to sign up for a later foreign-news digest sent after the study? & Binary yes / no choice & Zero-price later demand \\
\addlinespace
Any comments about the earlier materials or the later digest offer? & Open text & Descriptive feedback \\
\bottomrule
\end{tabularx}
\begin{tablenotes}[flushleft]
\footnotesize
\item Note: This table reports the non-matrix follow-up modules. Relative to endline, the follow-up survey was shorter and focused on persistence, retained learning, later demand, and perceived media bias.
\end{tablenotes}
\end{threeparttable}
\end{table}
